% GDZ PPN236018647,native552 with complete footnote continued on553;HKUST553–554.
\clearpage\SourceLeaf{552}\ArithmeticHead{523}{R}{1797 AUG. 1—1797 OCT.}
\NumberedItem{77}
\indent Demonstratum, viaque ad multa maiora per introd[uctionem] modulorum multiplicium strata.
\par\noindent\hfill[1797] Aug. 31.

\medskip\begingroup\small
\indent Vergl. die Nr. 76. Siehe den art. 372. der \textit{Analysis Residuorum}, Werke II, S. 237.
\par\noindent\hfill\textsc{Bachmann.}
\par\endgroup

\NumberedItem{78}
\indent Aug. 1. generalius ad quosvis modulos adaptatur.
\par\noindent\hfill[1797] Sept. 4.

\medskip\begingroup\small
\indent Hier ist Aug. 1. wohl ein Schreibfehler für Aug. 31., denn es handelt sich ohne Zweifel um den allgemeinsten Fall dessen, was in der Nr. 76 für einen Primzahlmodul, in der Nr. 77 für einen Modul, eine Primzahlpotenz ist, bewiesen worden und nun für einen ganz beliebigen Modul erledigt wird. Mit der Nr. 73 vom 1. August hat diese Notiz offenbar nichts zu tun.
\par\noindent\hfill\textsc{Bachmann.}
\par\endgroup

\NumberedItem{79}
\indent Principia detexi, ad quae congruentiarum secundum modulos multiplices resolutio ad congruentias secundum modulum linearem reducitur.
\par\noindent\hfill[1797] Sept. 9.

\medskip\begingroup\small
\indent Vergl. die Nr. 77 vom 31. August 1797, ferner die artt. 372 ff. der \textit{Analysis Residuorum}, Werke II, S. 237.
\par\noindent\hfill\textsc{Bachmann.}
\par\endgroup

\NumberedItem{80}
\indent Aequationes habere radices imaginarias methodo genuina demonstratum.
\par\noindent\hfill Bruns[vigae, 1797] Oct.
\par\indent Prom[ulgatum] in dissert[atione] pecul[iari\begingroup\renewcommand{\thefootnote}{*)}\footnotemark\endgroup] mense Aug. 1799.
\begingroup\renewcommand{\thefootnote}{[*)}\footnotetext{\textit{Demonstratio nova theorematis omnem functionem algebraicam rationalem integram unius variabilis in factores primi vel secundi gradus resolvi posse.} Helmstadii apud C. G. \textsc{Fleckeisen}, 1799.]}\endgroup

\medskip\begingroup\small
\indent Siehe Werke III, S. 1. \textsc{Gauss} ist auf Grund dieser Dissertation am 16. Juli 1799\begingroup\renewcommand{\thefootnote}{**)}\footnote{Siehe den Brief an W. \textsc{Bolyai} vom 16. Dezember 1799, \textit{Briefwechsel zwischen C. F. Gauss und W. Bolyai}, herausgegeben von F. \textsc{Schmidt} und P. \textsc{Stäckel}, Leipzig 1899, S. 34. In seiner Schrift: \textit{K. F. Gauss, Zwölf Kapitel aus seinem Leben}, Leipzig 1878, bezeichnet L. \textsc{Hänselmann}, S. 47, fälschlich den 14. Juli als den Tag, an dem »das Diplom ausgefertigt« wurde. Daß der 16. Juli der richtige Tag ist, wird bestätigt von P. \textsc{Zimmermann}, \textit{Zum Gedächtnis an Karl Friedrich Gauss}, Braunschweigisches Magazin \textbf{5}, 1899, Heft 15; dieser hat die Akten der Helmstedter Fakultät eingesehen. Das Gesuch von \textsc{Gauss} war am 26. Juni 1799 eingereicht worden und bereits am 28. Juni erstattete \textsc{Pfaff}, der die Dissertation schon kannte (vergl. die oben S. 99—105 abgedruckten Briefe), sein Gutachten. Die mündliche Prüfung wurde erlassen und das Diplom, wie es scheint, vor dem Druck der Dissertation ausgestellt, nachdem die Gebühren durch den Herzog von Braunschweig bezahlt worden waren.}\endgroup{} von der philosophischen Fakultät der Universität Helmstedt zum Doktor promoviert worden.
\par\noindent\hfill\textsc{Schlesinger.}
\par\endgroup
\par\vfill\noindent\hfill66\textsuperscript{*}
